\section{结论}
针对随机触发估计中的统计隐蔽攻击，本文研究了修改峰值与通信记录分布共同约束下的真实MSE损伤。
通过分解条件偏差平方与平均名义协方差变化，利用加权下一静默统计量建立损伤上界，
并构造满足完整联合日志约束的两拍因果攻击，得到阶数匹配的下界。
在全状态含噪测量、状态转移矩阵可逆的条件下，对于固定模型、有限维数、目标时刻及预设有限峰值范围，
包含下一事件的完整日志、事件串和精确协方差标记记录均对应 $\Theta(\min\{b^2,\varepsilon\})$ 的最坏损伤。

记录的观察时域和历史条件决定了上述保证的适用范围。
当拍截止日志可以保持正常分布而产生正损伤；标量反例进一步表明，删去历史条件后，下一事件的边缘分布也可能无法约束损伤。
二维耦合系统的仿真结果与损伤分解及双预算构造一致，并在同一攻击下展示了增加下一事件的作用。
本文给出了指定可信协议下的损伤阶数与记录条件，但上下界常数仍较保守，且依赖具体模型、维数和目标时刻。
后续可进一步研究常数的收紧，以及部分观测和时间一致情形下的损伤保证。
